% ============================================================
%  Chapter 3 — Memory and Storage
%  Inside a Smartphone: The Tiny World in Your Hand
% ============================================================

\chapter{Memory and Storage}

\chapteropening{%
  Your phone remembers your photos, your apps, and your games.
  But did you know it uses two completely different kinds of memory to do it?
}
\chapterbadge{Mission: Spot the Difference Between RAM and Storage!}

% --- Opening ---
\section*{Meet Memo}

\character{Memo}{Hi! I am Memo — and I have two jobs.
  My first job is RAM: holding things Chip is working on right now.
  My second job is storage: keeping your photos and apps safe for as long as you need them.
  They sound similar, but they are very different!}

\sectionrule

% --- RAM ---
\section*{RAM: The Working Memory}

\term{RAM} stands for \textbf{Random Access Memory}.
It is the phone's \emph{working memory} — the space where the processor
keeps information it is actively using.

\begin{picturethis}
  RAM is like a desk.
  When you are doing homework, you spread your books and papers across the desk.
  You can reach everything quickly.
  But when you finish and pack everything away, the desk is empty.

  RAM is the same: it holds things while the phone is working on them,
  but when you switch the phone off, everything on the ``desk'' disappears.
\end{picturethis}

When you open an app, it is loaded from storage into RAM so the processor
can work with it quickly. The more apps you have open at once,
the more RAM is being used — and the phone can slow down if RAM gets full.
When RAM is very full, the phone may close an app running in the background
to make space for the app you are using now.

\sectionrule

% --- Storage ---
\section*{Storage: The Long-Term Keep}

\term{Internal storage} is where the phone keeps everything permanently:
photos, videos, apps, messages, music, and the operating system itself.

Unlike RAM, storage \emph{does not} get wiped when the phone is turned off.
Your photos from last year are still there in the morning.

\begin{picturethis}
  Storage is like a bookshelf or filing cabinet.
  You can put hundreds of books and folders there, and they stay
  exactly where you left them — even overnight, even for years.
  To read a book, you take it off the shelf (load it into RAM) —
  the shelf does not empty just because you picked something up.
\end{picturethis}

Storage is measured in \term{gigabytes} (GB).
Most smartphones have between 64 GB and 512 GB of storage.

\begin{funfact}
  A single minute of 4K video can take up roughly 75 MB to 400 MB of storage,
  depending on quality settings.
  A 128 GB phone could hold about 5 to 28 hours of 4K footage — or around
  32,000 songs, or 40 million text messages.
\end{funfact}

\sectionrule

% --- Comparison ---
\section*{RAM vs Storage at a Glance}

\begin{quickmap}
  \begin{center}
  \begin{tabular}{lll}
    & \textbf{RAM} & \textbf{Storage} \\
    \hline
    \textbf{Purpose}  & Working memory & Long-term keeping \\
    \textbf{Speed}    & Very fast       & Slower than RAM \\
    \textbf{Cleared?} & Yes, when phone is off & No, stays forever \\
    \textbf{Typical size} & 4–12 GB      & 64–512 GB \\
    \textbf{Analogy}  & A desk         & A bookshelf \\
  \end{tabular}
  \end{center}
\end{quickmap}

\sectionrule

% --- Why phones run out of space ---
\section*{Why Phones Run Out of Space}

Photos, videos, apps, and downloaded music all take up storage space.
When storage gets full, the phone cannot save new things.

\begin{itemize}
  \item Delete photos and videos you no longer need.
  \item Remove apps you rarely use.
  \item Back up files to a computer or the cloud (storage on internet-connected computers).
\end{itemize}

\sectionrule

\begin{newwords}
  \begin{description}[labelwidth=3.5cm, leftmargin=3.7cm]
    \item[\term{RAM}]          Random Access Memory — fast, temporary working space.
    \item[\term{Storage}]      Permanent space for apps, photos, and files.
    \item[\term{Gigabyte (GB)}] A unit of data size: about 1 billion bytes.
    \item[\term{Operating system}] The main programme (iOS, Android) that manages the phone.
  \end{description}
\end{newwords}

\sectionrule

\begin{tryit}
  \textbf{Find out how much storage your device has.}

  With a grown-up's help, go to the phone's \emph{Settings} app
  and look for ``Storage'' or ``About Phone''.
  \begin{itemize}
    \item How much total storage does it have?
    \item How much is already used?
    \item What is using the most space — photos, apps, or something else?
  \end{itemize}
\end{tryit}

\sectionrule

\begin{bigidea}
  A smartphone uses two kinds of memory:
  \textbf{RAM} (fast, temporary working space cleared when off)
  and \textbf{storage} (permanent space for photos, apps, and files).
\end{bigidea}

\character{Memo}{Next, we will find out how you talk to the phone with your fingers.
  Meet Tappy in Chapter 4!}
